\documentclass[11pt,a4paper]{article}
\usepackage[margin=24mm]{geometry}
\usepackage{amsmath,amssymb,booktabs,array}
\usepackage[hidelinks]{hyperref}
\setlength{\parindent}{0pt}
\setlength{\parskip}{4pt}
\setlength{\emergencystretch}{2em}
\widowpenalty=10000
\clubpenalty=10000
\title{Continuous mass and ordered lanes\\for urban service-allocation gradients}
\author{Urban methodology review: consolidated draft}
\date{27 October 2026}
\begin{document}
\maketitle
\begin{abstract}
A vacant turn bay cannot serve traffic that cannot reach its opening.
We define a continuous-mass first-order traffic model that retains exact
route/movement order in realized local lanes. Known transfer capability
constrains progressive lane allocation; only admitted mass becomes
lane-committed. Unadmitted traffic retains priority within its movement,
while physical lane access determines whether other traffic can pass.
Ordered heads consume elapsed average service and share downstream
receiving among current competitors. The resulting timestep response is
composed into complete population time, including fixed clearance.
The scientific test is whether retained order improves externally
evaluated service decisions after strong macro models are allowed their
own conversion and inference from fine-grained observations.
\end{abstract}

\textbf{Scope.}
This consolidates the accepted urban task and closed traffic-policy
discussion into a model and test protocol. The companion specification
gives the full update and derivative branches. Implementation, native-rival
qualification and empirical usefulness remain separate work.

\section{The traffic question and its controls}

A through queue can obstruct the bay opening even when the bay has room.
More left service then need not help immediately, whereas through service
can remove the obstruction and later expose left demand. The same movement
totals can produce a different response when their realized order changes.
A macro model may still infer the relevant obstruction from its available
history; that success must count against the need for explicit order.

The accepted task is an isolated four-leg junction with fixed routes and
four protected effective-service groups: NS through/right, NS left,
EW through/right and EW left. The action satisfies
\begin{equation}
 g_p\geq0,\qquad \sum_{p=1}^4g_p=1,\qquad
 \mu_{\ell m}(g)=\min\{C_\ell,s_{\ell m}g_{p(m)}\}.
 \label{eq:compact-action}
\end{equation}
Fractions allocate average service, not executable phase intervals.
Traffic quantity remains continuous; route labels, permissions and topology
remain structural. A common loading period starts from empty. After the
control/release interval, a fixed candidate-independent policy clears
the complete population.

\begin{center}
\begin{tabular}{@{}p{0.08\linewidth}p{0.84\linewidth}@{}}
\toprule
V0 & Movement-exclusive paths: matched neutral propagation/service.\\
V1 & Shared lane: order-dependent head service without a bay.\\
V2 & Finite bay: storage case constructed without order-sensitive access.\\
V3 & Same bay layout: ordered access and storage interact.\\
V4 & Predeclared held-out variations of the admitted families.\\
\bottomrule
\end{tabular}
\end{center}

The starting bay configuration has a left-only bay and an adjacent shared
left/through/right lane. Other scenario permissions may be valid, but
V2/V3 pairs and compared models use the same declared matrix. Shared
permissions give left traffic a real queue-responsive choice; unique
permissions force its target.

\section{Allocation order and physical lane order}

Retain an upstream aggregate section $U$, a physical common feeder $F$
inside the local zone, then the bay $B$ and shared continuation $T$.
The fork is inside the zone. These are disjoint road pieces with their
own lengths, storage and histories. The mass-rank grid records movement
composition; rank does not locate a queue in metres or reconstruct exact
order from mixed bins. Ordered ledgers retain route references/positions
and continuous mass. Internal section crossings preserve route position;
the junction crossing advances it.

Two meanings of order must be kept distinct. An unassigned input record
sets the allocator's processing order. Once traffic occupies a real local
lane, its sequence determines physical front blocking. The allocator
may inspect a later movement after an unsuccessful allocation, but inspection
does not itself create a path past an occupied lane head.

For each ready homogeneous cohort, let $q$ be mass, $s$ remaining sending,
$u_\ell$ reachable legal receiving and $h_\ell$ the current queue-load
estimate. Feasible additions are selected together:
\begin{equation}
 V=\min(q,s,\sum_\ell u_\ell),\qquad
 x_\ell=\min\{u_\ell,[\lambda-h_\ell]_+\},\quad
 \sum_\ell x_\ell=V .
 \label{eq:compact-allocation}
\end{equation}
This fills the shortest legal queues subject to actual capability.
Only $x$ is committed. Later cohorts see $h+x$; remaining sending and
receiving are reduced by actual transfer. The unadmitted $q-V$ stays
with its old owner without a future-lane target. It retains priority
over later cohorts of its movement next update. Other movements are
still inspected. No frozen waiting-quota or old-target rationing
construction remains in the current model.

Queue perception uses pre-state local mass plus earlier actual admissions.
It does not invent lane loads for unassigned traffic or use forecast
signal plans as driver knowledge. In a compatible hand case with loads
$1.2/1.0$, available receiving $0.8/1.0$ and common sending $0.6$,
a ready left unit adds $0.2/0.4$, leaving balanced perceived loads
$1.4/1.4$ and unassigned upstream mass $0.4$.

\textbf{Physical access gives the scan its traffic meaning.}
Where paths are independently accessible, a blocked left movement does not
close an unrelated through entrance. Where traffic already occupies $F$,
a partly unserved front prevents later residents from reaching the fork.
The actual $U$--$F$--$B/T$ graph has no direct shortcut around $F$.
A full $F$ also restricts every movement using that shared space.

For example, with ready $F$ order $[T(0.4),L(0.6)]$, through receiving
$0.2$ and bay receiving $1$, only $0.2$ through enters; its residual
obstructs left access. Reverse the order with identical totals and an
initially shorter bay queue: left can use $0.6$ of bay room before
through uses $0.2$. The different totals, $0.2$ and $0.8$, follow from
physical order, rather than from globally freezing the logical allocator.
Both behaviours preserve the agreed processing policy and realized-lane FIFO.

A left turner's zero junction-service rate is not automatically a ban
on entering an accessible lane queue. A full bay is also not a ban on
using the legal shared lane if it is reachable. Admission space,
physical access and stop-line service refer to different constraints.

\section{Elapsed average service and causal propagation}

Each realized stop-line lane serves its current ready head at its
movement's average rate. Head clearance exposes the next resident.
A zero-rate head blocks that lane; other accessible lanes may continue.

At each head state, share downstream receiving green-weighted once per
inbound leg, then equally among its ready lanes subject to caps and
reuse. Newly ready heads join current competition. Earlier discharge
is not compensated to equalize final interval totals. The receiving
allowance for the update is represented by the selected constant rate
$R_j/\Delta t$. Head clearance remains sequential even though this avoids
a separately exhausted receiving pool.

With ready left mass $q_L$ ahead of sufficient through mass and
nonbinding receiving,
\begin{equation}
 y_T=s_Tg_T\left(\Delta t-\frac{q_L}{s_Lg_L}\right),
 \qquad 0<\frac{q_L}{s_Lg_L}<\Delta t.
 \label{eq:compact-service}
\end{equation}
Unit timestep/saturations, $q_L=0.2$ and fractions $0.4/0.6$ give
through discharge $0.3$. A small service transfer from through to left
has through-discharge derivative $0.25$ on this branch: the earlier
left-head clearance outweighs the reduced through rate. This is a
local response, not yet a complete-time benefit.

Propagation retains the repository's cumulative first-order sending and
receiving with fixed-delay, affine boundary-history interpolation.
Physical section lengths and delays remain compatible with the timestep.
Every node reads pre-state limits; today's departure does not create
immediate extra receiving, and new inflow cannot traverse its receiving
section in the same update. Internal head events calculate timestep
totals without becoming new propagation-history knots.

An internal transfer removes mass from one owner and adds the same mass
to the next, preserving route identity. Sending and receiving limits
keep stocks nonnegative and within storage. Unadmitted source or link
mass is counted once.

\section{One-step response and complete population time}

The complete objective includes source waiting, motion and queues until
all released demand leaves. Its inherited discrete account is
\begin{equation}
 J_{\Delta t}
 =\Delta t\sum_{k=0}^{H-1}
 \left[\sum_o(Q_{o,k}+D_{o,k})+\sum_en_{e,k}\right].
 \label{eq:compact-tst}
\end{equation}
At fixed starting state the immediate direct cost partial can be zero.
Changed service still affects future queues, access, spillback and sink
exit. An internal transfer alone does not remove population.

The specification differentiates the actual capped allocation, residual
updates, sharing, eligibility cuts, entrance reconstruction and head
durations on fixed active signatures. Compose those pieces as
\begin{equation}
 X_{k+1}=F_k(X_k,u_k),\qquad
 P_{k+1}=F_{k,X}P_k+F_{k,u}\dot u_k .
 \label{eq:compact-chain}
\end{equation}
The fixed continuation has $\dot u_k=0$, but existing state sensitivity
keeps propagating through clearance. There is no per-vehicle
trajectory-gradient construction.

For fixed releases, total population equals cumulative release minus
sink exit; the corresponding sink-exit derivative is an independent
check of complete-time sensitivity. Use a common horizon qualified to
clear the action neighbourhood, not a changing integer stop counter or
an unreported positive residual cutoff. Capacity/head/access switches
may create kinks; differentiability is conditional, not global.

\section{Strong comparisons and independent usefulness}

Macro rivals may use the original fine time series through their own
declared converter or estimator. Disclose retained state, memory,
observation schedule and calibration. Successful macro reconstruction
counts; exact ordered metadata retained internally remain meso information.
Equal current totals alone do not establish equality of lawful histories.
Separate state-conditioned replay, observation-updated prediction and
forecasting, stating what is fixed or regenerated for each derivative.

B0 is aggregate; B1 retains structured movement queues; B2 has responsive
lane/storage aggregates; B3 is the selective model. These roles are not
a universally nested information ladder. Native selection prioritizes
soundness/adaptation, task fit and then implementation ease.

The current B1 working source is
\href{https://arxiv.org/html/2310.12249v2}{Wei et al.'s link queue model}.
Its native effective-green and moving/queued formulation still require
a causal instantiated link/node treatment. B2's working direction uses
\href{https://www.sciencedirect.com/science/article/pii/S0191261525000578}
{Gong et al.'s lane-exit allocation} with an explicitly qualified storage
adaptation, for which
\href{https://attap.umd.edu/wp-content/uploads/2022/05/TRC522.pdf}
{Liu--Chang} is a candidate. Exit allocation does not itself define
finite-bay entry. Exact native implementation qualification remains
later work under the agreed selection policy.
The old exclusive-target mixture ablation is archived, not substituted
for a matched shared-permission comparator.

Independent hand/reference checks cover conservation, prefix access,
movement-specific scanning and receiving reuse. V0 checks neutral
equivalence under matched service and receiving. V2 must validate a
storage-aware explanation before V3 differences can support an order
claim. Rotation and properly symmetric lane limits check the physical
construction; equal lane totals alone do not justify one larger FIFO lane.

A later richer evaluator uses one declared fixed mapping from effective
fractions to executable signals, common routes/releases, fixed lane
behaviour/calibration and matched seeds. It must count the entire master
population from desired release to arrival, including insertion delay.
Unfinished traffic is retained or reported as noncompletion, not removed
by a completed-trip-only statistic.

First compare prescribed feasible service perturbations and their
externally realised objective changes. Only after forward and derivative
checks should matched optimization and held-out usefulness be tested.
Finite gradients or a decreasing internal objective are not that evidence.

\section{Deliverable and evidence boundary}

This consolidation applies the agreements, removes obsolete waiting
state, resolves processing versus physical order, and updates the
conditional sensitivity. The supporting specification supplies the
reconstructible update, feasibility reasoning and source/adaptation
index. It remains a written methodology for review; it does not claim
implementation correctness or empirical gradient usefulness.

\href{urban_methodology_specification.pdf}{Detailed mathematical specification}
and \href{urban_methodology_specification.tex}{retained LaTeX source}
are adjacent. Source-native adapters, numerical calibration, derivative
experiments and external evaluation remain qualified future work.
The previous question appendix and superseded formulas are preserved
in the dated archive, rather than interleaved with the current model.
\end{document}
